% Part: first-order-logic
% Chapter: natural-deduction
% Section: soundness

\documentclass[../../../include/open-logic-section]{subfiles}

\begin{document}

\iftag{FOL}
      {\olfileid[id]{fol}{ntd}{sou}}
      {\olfileid[id]{pl}{ntd}{sou}}

\olsection{Kesahihan}

\begin{explain}
Sebuah sistem !!{derivation}, seperti deduksi alami, disebut \emph{sahih} jika
sistem tersebut tidak dapat !!{derive} kesimpulan yang sebenarnya tidak
mengikuti premis-premisnya. Dengan demikian, kesahihan merupakan sejenis sifat keselamatan
yang dijamin bagi sistem !!{derivation}. Bergantung pada sifat teori-bukti
yang sedang dibahas, kita ingin mengetahui, misalnya, bahwa
\begin{enumerate}
\item setiap !!{sentence} yang !!{derivable} bersifat
  \iftag{FOL}{valid}{tautologis};
\item jika !!a{sentence} !!{derivable} dari beberapa kalimat lain,
  kalimat itu juga merupakan konsekuensi dari kalimat-kalimat tersebut;
\item jika suatu himpunan !!{sentence} tak konsisten, himpunan itu tak
  dapat dipenuhi.
\end{enumerate}
Ini adalah sifat-sifat penting dari suatu sistem !!{derivation}. Jika salah
satunya tidak berlaku, sistem !!{derivation} tersebut cacat---sistem itu
akan !!{derive} terlalu banyak. Oleh karena itu, membuktikan kesahihan
suatu sistem !!{derivation} sangatlah penting.
\end{explain}

\begin{thm}[Kesahihan]
\ollabel{thm:soundness}
Jika $!A$ !!{derivable} dari asumsi !!{undischarged}~$\Gamma$, maka
$\Gamma \Entails !A$.
\end{thm}

\begin{proof}
Misalkan $\delta$ adalah !!a{derivation} untuk $!A$. Kita menggunakan
induksi pada banyaknya inferensi dalam~$\delta$.

Untuk basis induksi, kita menunjukkan klaim tersebut jika banyaknya
inferensi adalah~$0$. Dalam kasus ini, $\delta$ hanya terdiri atas satu
!!{sentence}~$!A$, yakni sebuah asumsi. Asumsi itu !!{undischarged}, sebab
asumsi hanya dapat !!{discharged} oleh inferensi, sedangkan tidak terdapat
inferensi. Jadi, setiap
\iftag{FOL}{!!{structure}~$\Struct{M}$}{!!{valuation}~$\pAssign{v}$}
yang memenuhi semua asumsi !!{undischarged} dalam pembuktian tersebut
juga memenuhi~$!A$.

Sekarang untuk langkah induksi. Andaikan $\delta$ memuat~$n$ inferensi.
Premis-premis inferensi paling bawah kita !!{derive} melalui
sub-!!{derivation} yang masing-masing memuat kurang dari~$n$ inferensi.
Kita mengandaikan hipotesis induksi:
premis-premis inferensi paling bawah merupakan konsekuensi dari asumsi
!!{undischarged} dalam sub-!!{derivation} yang berakhir pada
premis-premis tersebut. Kita harus menunjukkan bahwa kesimpulan~$!A$
merupakan konsekuensi dari asumsi !!{undischarged} dalam seluruh
pembuktian.

Kita membedakan kasus menurut jenis inferensi paling bawah. Pertama, kita
membahas inferensi-inferensi yang mungkin dengan hanya satu premis.

\begin{enumerate}
\item Andaikan inferensi terakhir adalah \Intro{\lnot}:
  !!^{derivation} tersebut berbentuk
  \begin{prooftree}
    \AxiomC{$\Gamma, \Discharge{!A}{n}$}
    \RightLabel{$\delta_1$}
    \DeduceC{$\lfalse$}
    \DischargeRule{\Intro{\lnot}}{n}
    \UnaryInfC{$\lnot !A$}
  \end{prooftree}
  Menurut hipotesis induksi, $\lfalse$ merupakan konsekuensi dari asumsi
  !!{undischarged} $\Gamma \cup \{!A\}$ dalam~$\delta_1$. Perhatikan
  \iftag{FOL}{!!a{structure}~$\Struct{M}$}{!!a{valuation}~$\pAssign{v}$}.
  Kita harus menunjukkan bahwa jika
  $\iftag{FOL}{\Sat{M}}{\pSat{v}}{\Gamma}$, maka
  $\iftag{FOL}{\Sat{M}}{\pSat{v}}{\lnot !A}$. Untuk memperoleh
  kontradiksi, andaikan $\iftag{FOL}{\Sat{M}}{\pSat{v}}{\Gamma}$, tetapi
  $\iftag{FOL}{\Sat/{M}}{\pSat/{v}}{\lnot !A}$, yakni
  $\iftag{FOL}{\Sat{M}}{\pSat{v}}{!A}$. Ini berarti bahwa
  $\iftag{FOL}{\Sat{M}}{\pSat{v}}{\Gamma \cup \{!A\}}$, bertentangan
  dengan hipotesis induksi. Jadi,
  $\iftag{FOL}{\Sat{M}}{\pSat{v}}{\lnot !A}$.

\item Inferensi terakhir adalah \Elim{\land}. Ada dua varian: $!A$ atau
  $!B$ dapat diinferensikan dari premis $!A \land !B$. Perhatikan kasus
  pertama. !!^{derivation}~$\delta$ berbentuk sebagai berikut:
  \begin{prooftree}
    \AxiomC{$\Gamma$}
    \RightLabel{$\delta_1$}
    \DeduceC{$!A \land !B$}
    \RightLabel{\Elim{\land}}
    \UnaryInfC{$!A$}
  \end{prooftree}
  Menurut hipotesis induksi, $!A \land !B$ merupakan konsekuensi dari
  asumsi !!{undischarged}~$\Gamma$ dalam~$\delta_1$. Perhatikan
  \iftag{FOL}{!!a{structure}~$\Struct{M}$}{!!a{valuation}~$\pAssign{v}$}.
  Kita harus menunjukkan bahwa jika
  $\iftag{FOL}{\Sat{M}}{\pSat{v}}{\Gamma}$, maka
  $\iftag{FOL}{\Sat{M}}{\pSat{v}}{!A}$. Andaikan
  $\iftag{FOL}{\Sat{M}}{\pSat{v}}{\Gamma}$. Menurut hipotesis induksi
  ($\Gamma \Entails !A \land !B$), kita mengetahui bahwa
  $\iftag{FOL}{\Sat{M}}{\pSat{v}}{!A \land !B}$. Menurut definisi,
  $\iftag{FOL}{\Sat{M}}{\pSat{v}}{!A \land !B}$ jika dan hanya jika
  $\iftag{FOL}{\Sat{M}}{\pSat{v}}{!A}$ dan
  $\iftag{FOL}{\Sat{M}}{\pSat{v}}{!B}$. (Kasus ketika $!B$
  diinferensikan dari $!A \land !B$ ditangani dengan cara yang sama.)

\item Inferensi terakhir adalah \Intro{\lor}. Ada dua varian: $!A \lor
  !B$ dapat diinferensikan dari premis~$!A$ atau premis~$!B$. Perhatikan
  kasus pertama. !!^{derivation} tersebut berbentuk
  \begin{prooftree}
    \AxiomC{$\Gamma$}
    \RightLabel{$\delta_1$}
    \DeduceC{$!A$}
    \RightLabel{\Intro{\lor}}
    \UnaryInfC{$!A \lor !B$}
  \end{prooftree}
  Menurut hipotesis induksi, $!A$ merupakan konsekuensi dari asumsi
  !!{undischarged}~$\Gamma$ dalam~$\delta_1$. Perhatikan
  \iftag{FOL}{!!a{structure}~$\Struct{M}$}{!!a{valuation}~$\pAssign{v}$}.
  Kita harus menunjukkan bahwa jika
  $\iftag{FOL}{\Sat{M}}{\pSat{v}}{\Gamma}$, maka
  $\iftag{FOL}{\Sat{M}}{\pSat{v}}{!A \lor !B}$. Andaikan
  $\iftag{FOL}{\Sat{M}}{\pSat{v}}{\Gamma}$; maka
  $\iftag{FOL}{\Sat{M}}{\pSat{v}}{!A}$ karena $\Gamma \Entails !A$
  (hipotesis induksi). Jadi, harus pula berlaku bahwa
  $\iftag{FOL}{\Sat{M}}{\pSat{v}}{!A \lor !B}$. (Kasus ketika
  $!A \lor !B$ diinferensikan dari~$!B$ ditangani dengan cara yang sama.)

\item Inferensi terakhir adalah \Intro{\lif}. $!A \lif !B$
  diinferensikan dari subpembuktian dengan asumsi~$!A$ dan
  kesimpulan~$!B$, yakni
  \begin{prooftree}
    \AxiomC{$\Gamma, \Discharge{!A}{n}$}
    \RightLabel{$\delta_1$}
    \DeduceC{$!B$}
    \DischargeRule{\Intro{\lif}}{n}
    \UnaryInfC{$!A \lif !B$}
  \end{prooftree}
  Menurut hipotesis induksi, $!B$ merupakan konsekuensi dari asumsi
  !!{undischarged} dalam~$\delta_1$, yakni $\Gamma \cup \{!A\} \Entails
  !B$. Perhatikan
  \iftag{FOL}{!!a{structure}~$\Struct{M}$}{!!a{valuation}~$\pAssign{v}$}.
  Asumsi !!{undischarged} dalam~$\delta$ hanyalah $\Gamma$, sebab $!A$
  dilepaskan pada inferensi terakhir. Jadi, kita harus menunjukkan bahwa
  $\Gamma \Entails !A \lif !B$. Untuk memperoleh kontradiksi, andaikan
  bahwa untuk suatu
  \iftag{FOL}{!!{structure}~$\Struct{M}$}{!!{valuation}~$\pAssign{v}$},
  berlaku $\iftag{FOL}{\Sat{M}}{\pSat{v}}{\Gamma}$ tetapi
  $\iftag{FOL}{\Sat/{M}}{\pSat/{v}}{!A \lif !B}$. Jadi,
  $\iftag{FOL}{\Sat{M}}{\pSat{v}}{!A}$ dan
  $\iftag{FOL}{\Sat/{M}}{\pSat/{v}}{!B}$. Namun, menurut hipotesis, $!B$
  merupakan konsekuensi dari $\Gamma \cup \{!A\}$, yakni
  $\iftag{FOL}{\Sat{M}}{\pSat{v}}{!B}$, suatu kontradiksi. Jadi,
  $\Gamma \Entails !A \lif !B$.

\item Inferensi terakhir adalah \FalseInt. Di sini, $\delta$ berakhir
  dengan
  \begin{prooftree}
    \AxiomC{$\Gamma$}
    \RightLabel{$\delta_1$}
    \DeduceC{$\lfalse$}
    \RightLabel{\FalseInt}
    \UnaryInfC{$!A$}
  \end{prooftree}
  Menurut hipotesis induksi, $\Gamma \Entails \lfalse$. Kita harus
  menunjukkan bahwa $\Gamma \Entails !A$. Andaikan tidak; maka untuk
  suatu~$\iftag{FOL}{\Struct{M}}{\pAssign{v}}$ berlaku
    $\iftag{FOL}{\Sat{M}}{\pSat{v}}{\Gamma}$ dan
    $\iftag{FOL}{\Sat/{M}}{\pSat/{v}}{!A}$. Namun, selalu berlaku
    $\iftag{FOL}{\Sat/{M}}{\pSat/{v}}{\lfalse}$, sehingga ini berarti
    bahwa $\Gamma \Entails/ \lfalse$, bertentangan dengan hipotesis
    induksi.

\item Inferensi terakhir adalah \FalseCl: Latihan.

\iftag{FOL}{
\item Inferensi terakhir adalah \Intro{\lforall}. Maka $\delta$ berbentuk
  \begin{prooftree}
    \AxiomC{$\Gamma$}
    \RightLabel{$\delta_1$}
    \DeduceC{$!A(a)$}
    \RightLabel{\Intro{\lforall}}
    \UnaryInfC{$\lforall[x][!A(x)]$}
  \end{prooftree}
  Menurut hipotesis induksi, premis $!A(a)$ merupakan konsekuensi dari
  asumsi !!{undischarged}~$\Gamma$. Perhatikan suatu
  struktur~$\Struct{M}$ sebarang sedemikian sehingga $\Sat{M}{\Gamma}$. Kita harus
  menunjukkan bahwa $\Sat{M}{\lforall[x][!A(x)]}$. Karena
  $\lforall[x][!A(x)]$ adalah !!a{sentence}, ini berarti bahwa untuk setiap
  penugasan variabel~$s$, kita harus menunjukkan $\Sat{M}{!A(x)}[s]$
  (\olref[syn][ass]{prop:sat-quant}). Karena $\Gamma$ seluruhnya terdiri
  atas kalimat, $\Sat{M}{!B}[s]$ untuk semua $!B \in \Gamma$ menurut
  \olref[syn][sat]{defn:satisfaction}. Misalkan $\Struct{M'}$ sama dengan
  $\Struct{M}$ kecuali $\Assign{a}{M'} = s(x)$. Karena $a$ tidak muncul
  dalam~$\Gamma$, $\Sat{M'}{\Gamma}$ menurut
  \olref[syn][ext]{cor:extensionality-sent}. Karena $\Gamma \Entails
  !A(a)$, maka $\Sat{M'}{!A(a)}$. Karena $!A(a)$ adalah !!a{sentence},
  $\Sat{M'}{!A(a)}[s]$ menurut
  \olref[syn][ass]{prop:sentence-sat-true}. $\Sat{M'}{!A(x)}[s]$ jika dan
  hanya jika $\Sat{M'}{!A(a)}$ menurut
  \olref[syn][ext]{prop:ext-formulas} (ingat bahwa $!A(a)$ hanyalah
  $\Subst{!A(x)}{a}{x}$). Jadi, $\Sat{M'}{!A(x)}[s]$. Karena $a$ tidak
  muncul dalam~$!A(x)$, menurut \olref[syn][ext]{prop:extensionality},
  $\Sat{M}{!A(x)}[s]$. Namun, $s$ adalah penugasan variabel sebarang,
  sehingga $\Sat{M}{\lforall[x][!A(x)]}$.

\item Inferensi terakhir adalah \Intro{\lexists}: Latihan.

\item Inferensi terakhir adalah \Elim{\lforall}: Latihan.
}{}
\end{enumerate}

Sekarang mari kita bahas inferensi-inferensi yang mungkin dengan beberapa
premis: \Elim{\lor}, \Intro{\land}, \Elim{\lnot},
\iftag{FOL}{\Elim{\lif}, dan \Elim{\lexists}}{dan \Elim{\lif}}.
\begin{enumerate}

\item Inferensi terakhir adalah \Intro{\land}. $!A \land !B$
  diinferensikan dari premis $!A$ dan $!B$, dan $\delta$ berbentuk
  \begin{prooftree}
    \AxiomC{$\Gamma_1$}
    \RightLabel{$\delta_1$}
    \DeduceC{$!A$}
    \AxiomC{$\Gamma_2$}
    \RightLabel{$\delta_2$}
    \DeduceC{$!B$}
    \RightLabel{\Intro{\land}}
    \BinaryInfC{$!A \land !B$}
  \end{prooftree}
  Menurut hipotesis induksi, $!A$ merupakan konsekuensi dari asumsi
  !!{undischarged}~$\Gamma_1$ dalam~$\delta_1$, dan $!B$ merupakan
  konsekuensi dari asumsi !!{undischarged}~$\Gamma_2$ dalam~$\delta_2$.
  Asumsi !!{undischarged} dalam~$\delta$ adalah $\Gamma_1 \cup
  \Gamma_2$, sehingga kita harus menunjukkan bahwa $\Gamma_1 \cup
  \Gamma_2 \Entails !A \land !B$. Perhatikan
  \iftag{FOL}{!!a{structure}~$\Struct{M}$}{!!a{valuation}~$\pAssign{v}$}
  dengan $\iftag{FOL}{\Sat{M}}{\pSat{v}}{\Gamma_1 \cup \Gamma_2}$.
  Karena $\iftag{FOL}{\Sat{M}}{\pSat{v}}{\Gamma_1}$, harus berlaku
  $\iftag{FOL}{\Sat{M}}{\pSat{v}}{!A}$ karena $\Gamma_1 \Entails !A$;
  dan karena $\iftag{FOL}{\Sat{M}}{\pSat{v}}{\Gamma_2}$, berlaku
  $\iftag{FOL}{\Sat{M}}{\pSat{v}}{!B}$ karena $\Gamma_2 \Entails !B$.
  Bersama-sama, $\iftag{FOL}{\Sat{M}}{\pSat{v}}{!A \land !B}$.

\item Inferensi terakhir adalah \Elim{\lor}: Latihan.

\item Inferensi terakhir adalah \Elim{\lif}. $!B$ diinferensikan dari
  premis $!A \lif !B$ dan~$!A$. !!^{derivation}~$\delta$ berbentuk
  sebagai berikut:
  \begin{prooftree}
    \AxiomC{$\Gamma_1$}
    \RightLabel{$\delta_1$}
    \DeduceC{$!A \lif !B$}
    \AxiomC{$\Gamma_2$}
    \RightLabel{$\delta_2$}
    \DeduceC{$!A$}
    \RightLabel{\Elim{\lif}}
    \BinaryInfC{$!B$}
  \end{prooftree}
  Menurut hipotesis induksi, $!A \lif !B$ merupakan konsekuensi dari
  asumsi !!{undischarged}~$\Gamma_1$ dalam~$\delta_1$, dan $!A$
  merupakan konsekuensi dari asumsi !!{undischarged}~$\Gamma_2$
  dalam~$\delta_2$. Perhatikan
  \iftag{FOL}{!!a{structure}~$\Struct{M}$}{!!a{valuation}~$\pAssign{v}$}.
  Kita harus menunjukkan bahwa jika
  $\iftag{FOL}{\Sat{M}}{\pSat{v}}{\Gamma_1 \cup \Gamma_2}$, maka
  $\iftag{FOL}{\Sat{M}}{\pSat{v}}{!B}$. Andaikan
  $\iftag{FOL}{\Sat{M}}{\pSat{v}}{\Gamma_1 \cup \Gamma_2}$. Karena
  $\Gamma_1 \Entails !A \lif !B$, maka
  $\iftag{FOL}{\Sat{M}}{\pSat{v}}{!A \lif !B}$. Karena
  $\Gamma_2 \Entails !A$, kita memperoleh
  $\iftag{FOL}{\Sat{M}}{\pSat{v}}{!A}$. Ini berarti bahwa
  $\iftag{FOL}{\Sat{M}}{\pSat{v}}{!B}$. (Sebab, jika
  $\iftag{FOL}{\Sat/{M}}{\pSat/{v}}{!B}$, karena
  $\iftag{FOL}{\Sat{M}}{\pSat{v}}{!A}$, akan berlaku
  $\iftag{FOL}{\Sat/{M}}{\pSat/{v}}{!A \lif !B}$, bertentangan dengan
  $\iftag{FOL}{\Sat{M}}{\pSat{v}}{!A \lif !B}$.)

\item Inferensi terakhir adalah \Elim{\lnot}: Latihan.

\tagitem{FOL}{Inferensi terakhir adalah \Elim{\lexists}: Latihan.}{}
\end{enumerate}
\end{proof}

\tagprob{FOL}
\begin{prob}
Lengkapi pembuktian \olref[fol][ntd][sou]{thm:soundness}.
\end{prob}
\tagendprob

\tagprob{notFOL}
\begin{prob}
Lengkapi pembuktian \olref[pl][ntd][sou]{thm:soundness}.
\end{prob}
\tagendprob

\begin{cor}
\ollabel{cor:weak-soundness}
Jika $\Proves !A$, maka $!A$ \iftag{FOL}{valid}{merupakan tautologi}.
\end{cor}

\begin{cor}
\ollabel{cor:consistency-soundness}
Jika $\Gamma$ dapat dipenuhi, maka himpunan tersebut konsisten.
\end{cor}

\begin{proof}
Kita membuktikan kontraposisinya. Andaikan $\Gamma$ tidak konsisten.
Maka $\Gamma \Proves \lfalse$, yakni terdapat !!a{derivation} untuk
$\lfalse$ dari asumsi !!{undischarged} dalam~$\Gamma$. Menurut
\olref{thm:soundness}, setiap
\iftag{FOL}{!!{structure}~$\Struct{M}$}{!!{valuation}~$\pAssign{v}$}
yang memenuhi $\Gamma$ harus memenuhi $\lfalse$. Karena
$\iftag{FOL}{\Sat/{M}}{\pSat/{v}}{\lfalse}$ untuk setiap
\iftag{FOL}{!!{structure}~$\Struct{M}$}{!!{valuation}~$\pAssign{v}$},
tidak ada \iftag{FOL}{$\Struct{M}$}{$\pAssign{v}$} yang dapat memenuhi
$\Gamma$; yakni, $\Gamma$ tidak dapat dipenuhi.
\end{proof}

\end{document}
